\section{Live Demos y límites de la acción}

En clase se usan la compra de boletos de avión, la operación en dispositivos móviles, un examen de manejo y la Action API para mostrar la capacidad de acción visible de un agent. Una lectura de calidad no puede limitarse a repetir «lo logró»; debe preguntarse paso a paso: cómo se introduce el objetivo, qué devuelve el entorno, dónde ocurre un cambio de estado, qué pasos requieren confirmación y cómo se detiene o se recupera el sistema cuando falla.

\subsection{Demo de compra de boletos de avión: de la solicitud al itinerary}

El primer grupo de cinco cuadros registra un flujo de browser-agent flight-booking. La página inicial presenta primero el contexto de la tarea; después, la interfaz de desarrollo envía la solicitud, el navegador pasa a los resultados de búsqueda y al checkout, y al final muestra el itinerary. La secuencia de estados indica que el agent puede sostener la tarea a lo largo de varias páginas web, pero la grabación no informa cuántas veces se repitió la prueba, cómo se distribuyen los fallos ni todas las llamadas en segundo plano.

\lecturefigure{slide-32-flight-demo-start.jpg}{Estado inicial de la demo de compra de boletos de avión}{Video de clase de Stanford 00:23:30}

\begin{knowledgebox}{Lectura de la figura: confirmar primero los límites de la tarea, no la velocidad de la animación}
El cuadro inicial solo demuestra que la demostración está por entrar en la tarea de vuelos. El lector debe registrar si la cuenta de partida, la ciudad de destino, las fechas, la preferencia de aerolínea y los permisos de pago estaban configurados de antemano; este hidden state influye de forma notable en la dificultad de la tarea y también determina si el resultado puede reproducirse.
\end{knowledgebox}

\lecturefigure{slide-33-flight-demo-code.jpg}{Envío de la solicitud de compra al Agent mediante la interfaz de desarrollo}{Video de clase de Stanford 00:23:36}

\begin{knowledgebox}{Lectura de la figura: la instrucción debe convertirse en restricciones ejecutables}
Una solicitud en lenguaje natural suele contener restricciones duras y preferencias flexibles. El sistema debe marcar la fecha, los aeropuertos, el número de personas y el presupuesto como hard constraints, y la aerolínea United, la franja de salida o el asiento como soft preferences; si ambas entran en conflicto, no puede sacrificar en silencio una restricción dura.
\end{knowledgebox}

\lecturefigure{slide-34-flight-demo-results.jpg}{El Agent recorre los resultados de búsqueda de vuelos}{Video de clase de Stanford 00:24:33}

\begin{knowledgebox}{Lectura de la figura: la página de resultados es un problema de elección, no un estado de tarea completada}
Este cuadro indica que el entorno ya devolvió vuelos candidatos. El Agent todavía debe comparar precio, escalas, zonas horarias, reglas de equipaje y preferencias, y guardar un audit trail de «por qué eligió esta opción». Que aparezca la página no significa que el candidato ya satisfaga el objetivo del usuario.
\end{knowledgebox}

\lecturefigure{slide-35-flight-demo-checkout.jpg}{El Agent entra en la etapa de checkout}{Video de clase de Stanford 00:24:51}

\begin{warningbox}{El checkout es un límite de permisos}
Buscar y llenar formularios suele ser reversible; la compra final puede generar cargos y costos de cancelación. Al entrar al checkout se debe mostrar un resumen estructurado y pedir confirmación explícita antes de ejecutar la purchase; una tarjeta de crédito guardada no puede interpretarse como una autorización permanente para cualquier transacción.
\end{warningbox}

\lecturefigure{slide-36-flight-demo-itinerary.jpg}{Estado final del itinerary en la grabación de la clase}{Video de clase de Stanford 00:25:15}

\begin{knowledgebox}{Lectura de la figura: la página final todavía requiere verificar la postcondition}
El sistema debe cotejar pasajero, fecha, aeropuertos, número de vuelo, precio total y confirmation number, y escribir el resultado en un registro consultable. Si la UI solo muestra el itinerary sin que se verifiquen los campos clave, todavía puede producirse un falso éxito de «flujo completado, tarea equivocada».
\end{knowledgebox}

\teachervoice{El ponente afirma que la demo ya estaba personalizada con su preferencia por United, su cuenta y su método de pago, y subraya que el agent tiene purchasing power. La señal más valiosa de la clase no es la exhibición técnica, sino que vincula directamente personalization, credential access e irreversible action.}

\begin{lstlisting}[caption={Máquina de estados de compra de boletos con puerta de confirmación}]
state = collect_constraints(request, user_profile)
candidates = search_flights(state)
choice = rank_and_explain(candidates, state)
draft = fill_checkout(choice)
assert validate(draft, state.hard_constraints)
approval = ask_user_confirmation(draft.summary)
if approval:
    receipt = execute_purchase(draft)
    verify_postconditions(receipt, state)
else:
    stop_without_charge()
\end{lstlisting}

\subsection{Demo móvil: las restricciones del entorno interrumpen el plan}

El segundo grupo de cuadros lleva la misma idea de control por lenguaje natural al teléfono. La página principal de la platform, la solicitud por mensaje, la web search, el food result y el delivery-area failure forman una trayectoria con más valor didáctico, porque el último cuadro muestra explícitamente una limitación del entorno en lugar de recortarse como un video promocional de éxito total. Por eso esta sección toma como eje «cómo un fallo expone las restricciones del entorno», y no solo cuenta cuántas aplicaciones se recorrieron.

\lecturefigure{slide-37-mobile-platform.jpg}{Entrada a la plataforma móvil de MultiOn}{Video de clase de Stanford 00:25:21}

\begin{knowledgebox}{Lectura de la figura: el dispositivo móvil es una nueva superficie de I/O}
El Agent debe manejar una UI de pantalla pequeña, permisos del sistema, notificaciones, cambios entre apps y el estado de la red. La entrada en lenguaje natural unifica la expresión de la intención, pero no unifica el estado de las aplicaciones subyacentes; cada app puede seguir teniendo su propio inicio de sesión, ubicación geográfica y reglas de confirmación.
\end{knowledgebox}

\lecturefigure{slide-38-mobile-message.jpg}{Expresión de la tarea mediante una interfaz de mensajes}{Video de clase de Stanford 00:25:39}

\begin{knowledgebox}{Lectura de la figura: el mensaje es un objetivo, no una especificación completa}
Una solicitud breve suele omitir presupuesto, cantidad, dirección y alternativas. El sistema debe leerlos de un profile confiable o preguntar; no puede usar conjeturas del modelo de lenguaje para llenar campos que tienen consecuencias económicas o de privacidad.
\end{knowledgebox}

\lecturefigure{slide-39-mobile-web-search.jpg}{El Agent móvil lanza una web search}{Video de clase de Stanford 00:26:03}

\begin{knowledgebox}{Lectura de la figura: la ejecución entre aplicaciones requiere un estado de tarea compartido}
Al pasar del mensaje a la búsqueda, el objetivo, las restricciones y las fuentes deben mantenerse coherentes. Si cada herramienta solo recibe texto libre, la información tiende a desviarse en las sucesivas reformulaciones; un task state estructurado permite que la búsqueda, el filtrado y el pedido compartan el mismo conjunto de campos.
\end{knowledgebox}

\lecturefigure{slide-40-mobile-food-result.jpg}{El Agent encuentra resultados de food ordering}{Video de clase de Stanford 00:26:21}

\begin{knowledgebox}{Lectura de la figura: que un candidato sea visible no significa que pueda entregarse}
La página del restaurante o del producto solo indica que la búsqueda tuvo éxito. La postcondition real incluye que la dirección de entrega esté dentro del área de cobertura, que el producto esté disponible, que el costo sea aceptable y que el tiempo estimado cumpla lo requerido. Un Agent benchmark debe incluir estas restricciones del entorno en la success definition.
\end{knowledgebox}

\lecturefigure{slide-41-mobile-delivery-limit.jpg}{La restricción del área de entrega bloquea la tarea}{Video de clase de Stanford 00:26:30}

\begin{warningbox}{El fallo debe convertirse en un estado explícito}
Una delivery-area constraint no es un error de razonamiento que desaparezca porque el modelo «piense un poco más». La conducta correcta es informar la causa del bloqueo, ofrecer alternativas viables y esperar a que el usuario elija, en lugar de hacer clic en bucle, ocultar el fallo o cambiar la dirección por cuenta propia.
\end{warningbox}

\teachervoice{Lo que más vale la pena conservar de este grupo de demos es el cuadro del fallo: el entorno le dice al agent que el plan actual no es viable. Anticipa el tema posterior de la plan divergence: sin retroalimentación real, un plan coherente en el lenguaje se aleja cada vez más del mundo ejecutable.}

\subsection{Driving-test Claim e human-like interface}

A continuación, la clase muestra una diapositiva con un marcador de posición MP4, y el ponente afirma que el agent puede completar un driving test en línea. Como la grabación no presenta las preguntas completas, las respuestas, la calificación ni experimentos repetidos, las notas lo registran como una capability claim y no como un resultado verificado de forma independiente. La siguiente diapositiva explica, desde la idea de «manos y ojos digitales e interfaces de tipo humano», por qué los investigadores quieren que el agent use directamente el software existente.

\lecturefigure{slide-42-driving-test.jpg}{La online driving-test demo mencionada en clase}{Video de clase de Stanford 00:26:36}

\begin{warningbox}{Una Claim no sustituye a una evaluation}
Para verificar la capacidad en el driving-test se necesita, como mínimo, un conjunto público de preguntas o un protocolo de prueba cerrado, una revisión de filtración de preguntas, aleatorización, múltiples ejecuciones y un análisis de errores por categoría. Un ícono de archivo de video no puede sustentar conclusiones sobre tasa de aprobación, generalización o seguridad.
\end{warningbox}

\lecturefigure{slide-43-human-like-agents.jpg}{Why human-like AI agents: reutilizar entornos digitales diseñados para personas}{Video de clase de Stanford 00:28:27}

\begin{knowledgebox}{Lectura de la figura: human-like se refiere a compatibilidad de interfaz, no a equivalencia cognitiva}
El entorno digital ya cuenta con teclado, ratón, disposición visual y documentación en lenguaje natural. Que el agent use estas interfaces permite reutilizar el software existente y evita desarrollar primero una API para cada sitio web; el costo es la ambigüedad visual, la fragilidad de los clics, el estado oculto y una superficie de ataque mayor. Lo que describe es I/O compatibility, no que el agent entienda el mundo como una persona.
\end{knowledgebox}

\subsection{Cinco niveles de Autonomy: delegar permisos de forma gradual según el riesgo}

La Autonomy slide toma prestada la intuición de los niveles de la conducción autónoma y lleva al agent desde el control humano total hacia una alta autonomía. El texto de la figura es pequeño; el objetivo didáctico no debería ser memorizar los nombres de los niveles, sino entender qué cambia en cada nivel: quién propone el plan, quién ejecuta las acciones, quién supervisa y quién asume la responsabilidad de la recuperación.

\lecturefigure{slide-44-autonomy-levels.jpg}{Propuesta de cinco niveles de Agent Autonomy presentada en clase}{Video de clase de Stanford 00:28:45}

\begin{importantbox}{Diseñar la autonomy según el riesgo de la acción}
\begin{tabularx}{\textwidth}{L{0.10\textwidth}>{\raggedright\arraybackslash}X >{\raggedright\arraybackslash}X}
\toprule
Nivel & Papel del sistema & Control recomendado \\
\midrule
L1 & Solo sugiere; el usuario ejecuta & Mostrar fundamentos y pasos editables \\
L2 & Ejecuta acciones reversibles de bajo riesgo & Cada lote de acciones se puede previsualizar y cancelar \\
L3 & Maneja tareas de varios pasos con confirmación en puntos clave & Colocar un gate antes de dinero, mensajes salientes o cambios de cuenta \\
L4 & Autónomo dentro de un dominio acotado & Límites de presupuesto, dominios, tiempo y permisos \\
L5 & Alta autonomía en un dominio amplio & Visión de la clase; el despliegue real todavía requiere supervisión externa y kill switch \\
\bottomrule
\end{tabularx}
\end{importantbox}

Si la pérdida potencial de la acción $a$ es $L(a)$, su probabilidad de error es $q(a)$ y su coeficiente de recuperabilidad es $r(a)\in[0,1]$, puede usarse
\[
R(a)=q(a)L(a)(1-r(a))
\]
como ordenamiento aproximado del riesgo. Cuanto mayor sea $R(a)$, más deben reducirse los permisos, añadirse confirmaciones o prohibirse la ejecución automática; esta fórmula es una abstracción de ingeniería de las notas, no un estándar dado en clase.

\teachervoice{El ponente compara la autonomy con el proceso de niveles de la self-driving y reconoce que la full autonomy todavía enfrenta muchos problemas. La ruta sólida no es proclamar primero «totalmente automático», sino separar en capas las acciones reversibles de las irreversibles y ampliar gradualmente el dominio de autonomía ya verificado.}

\subsection{Disyuntiva entre API y direct interaction}

La Computer Interaction slide compara dos rutas: operar mediante la API que ofrece la aplicación, o hacer que el agent use directamente keyboard, mouse y la pantalla. La API suele ser estructurada, verificable y con permisos claros; la direct interaction tiene mayor cobertura, pero depende más de la estabilidad de la UI y de la localización visual.

\lecturefigure{slide-45-computer-interaction.jpg}{Las dos rutas de la Agent computer interaction}{Video de clase de Stanford 00:31:06}

\begin{knowledgebox}{Términos clave: API y operación directa}
\begin{tabularx}{\textwidth}{L{0.18\textwidth}X X}
\toprule
Ruta & Ventajas & Riesgos principales \\
\midrule
API interaction & Parámetros y valores de retorno estructurados; facilita schema validation, reintentos y mínimo privilegio & Cobertura limitada; requiere soporte del proveedor del servicio; la API puede cambiar \\
Keyboard/mouse & Reutiliza casi cualquier interfaz humana; despliegue rápido de prototipos & Coordenadas frágiles, ventanas emergentes ocultas, reconocimiento visual erróneo, prompt injection y clics equivocados \\
Hybrid & Prioriza la API, recurre a la UI cuando falta, y comparte el task state & Enrutamiento y consistencia más complejos; debe marcarse el canal de ejecución actual \\
\bottomrule
\end{tabularx}
\end{knowledgebox}

\teachervoice{La clase plantea explícitamente el trade-off: la API es más safe y controllable; la direct interaction, más general. Las notas lo convierten en un principio de ingeniería: la generalidad no es un beneficio gratuito; cuanto más libre es la interfaz, más importantes son la verificación del entorno y el aislamiento de permisos.}

\subsection{Action API: convertir el lenguaje natural en acciones controladas}

El último grupo de cinco cuadros muestra el overview de la Action API, la llamada desde código, la interfaz de chat, el proceso de control y el resultado. Busca ofrecer una capa de lenguaje natural para que los desarrolladores no tengan que escribir a mano una automatización específica para cada acción de UI; pero la palabra «Action» del nombre de la API no debe ocultar el flujo de control: la solicitud todavía debe analizarse, ejecutarse, observarse, verificarse y detenerse. Esta sección lee cada cuadro siguiendo esa cadena de control, para no confundir una entrada unificada con una confiabilidad unificada.

\lecturefigure{slide-46-action-api-overview.jpg}{Presentación en clase de la Action API de MultiOn}{Video de clase de Stanford 00:32:18}

\begin{knowledgebox}{Lectura de la figura: una entrada unificada no implica una semántica unificada}
La Action API puede unificar la forma de llamada, pero los sitios web, las aplicaciones y los permisos subyacentes siguen siendo distintos. Una implementación confiable debe convertir el lenguaje natural en una typed intent y conservar el dominio objetivo, las acciones permitidas, el número máximo de pasos y la condición de éxito.
\end{knowledgebox}

\lecturefigure{slide-47-action-api-code.jpg}{El desarrollador llama a la Action API desde código}{Video de clase de Stanford 00:32:21}

\begin{knowledgebox}{Lectura de la figura: la llamada al SDK es solo el punto de partida del plano de control}
Del lado del código se deben registrar, como mínimo, request id, policy, budget y callback; del lado de la ejecución se devuelven structured observations, y no solo un «listo». Así quien llama puede distinguir entre en ejecución, requiere confirmación, fallido y completado con verificación.
\end{knowledgebox}

\lecturefigure{slide-48-action-api-chat.jpg}{Action API y entrada de tareas en formato de chat}{Video de clase de Stanford 00:32:33}

\begin{knowledgebox}{Lectura de la figura: la capa de chat se encarga de aclarar ambigüedades}
Cuando a la solicitud le falta el objeto, el horario, el presupuesto o el destinatario, la chat interface debe preguntar por iniciativa propia. Enviar todo el lenguaje natural directamente al ejecutor convierte la conversational ambiguity en errores en el mundo real.
\end{knowledgebox}

\lecturefigure{slide-49-action-api-control.jpg}{El Agent toma el control de la interfaz de la aplicación para ejecutar acciones}{Video de clase de Stanford 00:32:42}

\begin{warningbox}{Al controlar la interfaz hay que evitar el desfase entre observación y acción}
El desplazamiento de la página, las ventanas emergentes o la carga asíncrona pueden invalidar coordenadas anteriores. Antes de cada acción debe volver a observarse el elemento objetivo, y después de ella debe comprobarse el cambio de estado esperado; si la observation no coincide, hay que detenerse en lugar de seguir con un plan caducado.
\end{warningbox}

\lecturefigure{slide-50-action-api-result.jpg}{Estado del resultado de la demo de Action API}{Video de clase de Stanford 00:33:06}

\begin{knowledgebox}{Lectura de la figura: el Result debe vincularse a un success predicate}
La pantalla de resultado solo cuenta como éxito si cumple un predicate declarado de antemano, por ejemplo «el mensaje objetivo se envió al destinatario correcto y con el contenido correcto». Que aparezca visualmente la app objetivo o el texto no basta para demostrar que el estado del backend ya se confirmó.
\end{knowledgebox}

\begin{lstlisting}[caption={Controlador observe-act-verify de la Action API}]
for step in range(max_steps):
    observation = observe_environment()
    decision = policy.plan(task, observation, permissions)
    if decision.requires_confirmation:
        decision = confirm_with_user(decision)
    result = execute(decision.action)
    if satisfies_success_predicate(result, task):
        return verified_success(result)
    if result.is_blocked or result.risk_exceeded:
        return safe_stop(result)
return timeout_without_claiming_success()
\end{lstlisting}

\subsection{Resumen de la sección}

Una Live demo puede mostrar la posibilidad end-to-end, pero por sí sola no aporta una tasa de éxito general. El flujo de compra de boletos revela la compra irreversible; el fallo en el móvil revela las restricciones del entorno; los niveles de autonomy revelan que la delegación de permisos debe avanzar según el riesgo, y la Action API muestra que una entrada unificada en lenguaje natural sigue necesitando typed state, confirmation gate y postcondition verification.
